% ===================================================================
%  Section XI: Conclusion and Future Work
% ===================================================================

\subsection{Summary of Contributions}

This paper presented a layered artifact system for zero-knowledge verifiable
offline DQN training fragments over provenance-bound committed replay data.
The key technical contributions are:

\begin{enumerate}[leftmargin=1.8em]
  \item A provenance-bound dataset pipeline that separates self-collected
  replay/reward audits from public benchmark source-integrity imports,
  binding all provenance hashes into a SHA-256 Merkle commitment.

  \item Three constraints the fragment relation declares and enforces beyond
  recomputing the update: a sampler seed derived from the dataset commitment
  rather than supplied by the prover, a published bound on the fixed-point
  values checked before the multiplication that overflows, and component-wise
  gradient clipping.

  \item SP1 proof-backed relations covering replay membership (with
  dataset-size scaling to 100k leaves), Bellman/TD arithmetic, forward MLP
  evidence, one-step SGD updates, checkpoint chains, multi-step fragments,
  chunk-chain aggregation over child manifests ($T \in \{32, 64, 128\}$), and
  recursive aggregation that verifies each child proof inside the guest
  ($T \in \{16, 32, 64\}$ flat, and binary trees up to $T = 4{,}992$).

  \item Single root proofs covering a complete training run: $1{,}248$ steps on
  both environments and $4{,}992$ steps on \texttt{lunarlander-random}, the
  last covering the run that produced a policy reported in
  \Cref{tab:rl-performance}.

  \item Ten formal security statements, each mapped to implemented relation
  code, SP1 backend, tests, and benchmark evidence.

  \item An evaluation that measures what the provable configuration costs in
  policy quality rather than asserting it is small, and that reports verify
  times of $0.054$--$0.197$\,s and proof sizes of $1.27$--$4.31$\,MB across
  twenty-four of the twenty-five proved rows, naming the Groth16 variant as the
  exception on both axes.
\end{enumerate}

The contribution is a deliberately scoped security/systems artifact.
Every formal theorem maps to a relation, backend, tests, and benchmark
evidence. The reviewer fast path (\texttt{make reproduce-small}) exercises
the entire pipeline without heavy proof regeneration.

\subsection{Future Work}

Several directions can extend the current system:

\textbf{Future work: recursion beyond $T = 64$ and off the GPU.} Recursive
child-proof verification inside the aggregate SP1 guest is implemented and
measured for $T \in \{16, 32, 64\}$, which removes the need for external
child-proof verification at those sizes. Two limits remain. Peak device memory
is flat in $T$, so longer traces are plausible but untested beyond $T = 64$.
And the CPU prover does not complete these configurations within $61$~GB, so
reproducing them presently requires a CUDA device with at least $20$~GB.

\textbf{Adam optimizer proofs.} Covering Adam would need its per-parameter
adaptive rates, first and second moment estimates, and bias correction inside
the relation. There is a prior obstacle that is easy to miss: at
$\fpscale = 1000$ the library default rate of $3\times 10^{-4}$ encodes to
zero, so Adam is not representable at its own default before any of the moment
arithmetic is considered. A finer fixed-point scale would be a prerequisite,
and would raise the overflow bound in turn.

\textbf{Batch-size scaling.} The relation proves batch-size-1 updates, and
supporting larger batches would need vectorized fixed-point gradient
accumulation at a cycle cost that grows with the batch. We had expected this to
be the change that closes the gap to a tuned implementation. The sweep in
\Cref{subsec:provable-config} says otherwise: at a long enough target-sync
interval, batch size 1 reaches $-152.7$ on \texttt{lunarlander-random} against
$-215.3$ for the tuned minibatch row. Batching is worth having, but it is not
what was holding the provable configuration back.

\textbf{Larger network architectures.} Moving beyond the canonical tiny
Linear--ReLU--Linear MLP to deeper networks with convolutional layers or
batch normalization would extend applicability to Atari-scale DQN.

\textbf{On-chain verification.} Compressing SP1 STARK proofs to
Groth16~\cite{groth2016zksnark} SNARKs would enable on-chain verification.
We measured the adjacent configuration -- verifying Groth16 child proofs
\emph{inside} the aggregate guest -- and it is expensive on both axes
(\Cref{tab:proof-cost}). Wrapping the final root proof for an on-chain verifier
is a different and more promising operation, which we have not measured.

\textbf{Honest data collection proofs.} Extending the provenance pipeline to
provide cryptographic evidence of honest data collection (e.g., through
trusted execution environments or verifiable random functions) would
strengthen the end-to-end trust chain.
